\section{Literature Review}
\subsection{Sound Source Localization}

The first part of the proposed project (before the motion control itself) involves deducing the target's position by sensing its acoustic transmissions. Sound Source Localization (SSL) is a mature field from which the project can borrow heavily. \citet{rascon2017review} and
\citet{argentieri2015survey} both conducted extensive reviews on SSL methods.

\citet{rascon2017review} mostly look at strategies used by arrays of omnidirectional microphones (usually 2) to locate one or multiple sound sources. The main strategy across papers for angle estimation is to process the received signal to extract certain statistics, use a grid search or a trained model to map these statistics to a location, then to smooth the result via different filters (depending on the specific setup). Distance estimation is more tricky: the extracted signal statistics correlates more weakly with source distance compared to source angle, especially when the source's baseline loudness is unknown or multiple sources are transmitting at once. Some surveyed works choose to leverage robot movement to resolve this (which is the approach most relevant to our setup), by triangulating using measurements taken from different positions, sometimes explicitly moving the robot closer to the source. 

\citet{argentieri2015survey} look separately at binaural methods (which use two
microphones to imitate human hearing) and array methods (more than 2 microphones). Binaural
systems rely on the interaural time and level differences (intensity is limited in usefulness when the source's loudness is unknown). Array methods use a variety of techniques such as beamforming (which shifts each microphone’s signal by the delays a candidate direction predicts and keeps the direction that results in the highest array energy), Generalized Cross-Correlation with Phase Transform (GCC-PHAT), or Multiple Signal Classification (MUSIC).
Array methods are found to be overall more robust, and the review stresses the under-utilization of robot motion in SSL (what is called \emph{active audition} \cite{nakadai2000active}). Several techniques can leverage active audition to refine estimates, for example by intersecting successive ``cones of confusion'' (the set of directions consistent with the measurements at each position).

The closest application in SSL is found in the work of
\citet{luo2010docking}. They built an indoor robot that finds a sound-emitting charging dock using two microphone pairs, a camera and laser scanner.

\subsection{Rotating Sensor Arrays and Leveraging Robot Motion}

The specific strategy of rotating a directional sensor and using its amplitude readings to deduce source location is well-established, although it is more commonly found in radio applications. \citet{graefenstein2009wireless} are an early example of the setup, using a rotating directional antenna on a ground robot to estimate the angle of arrival (AoA) of a wireless node from the detected signal strength during the rotation.

There is significant design freedom in how exactly each sweep becomes a direction estimate. There is no consensus, and methods range from simple heuristics to more involved estimates: \citet{graefenstein2009wireless}
take the angle of strongest signal (a natural choice); \citet{venkateswaran2013rf} split the angular space into bins and average the readings within each to extract the loudest bin; \citet{tokekar2011active} sample the signal strength at
fixed angles and fit a polynomial around the maximum. Some methods take advantage of the whole sweep rather than discarding most readings: \citet{singleantenna2025} model the antenna's known radiation pattern to derive a maximum-likelihood (ML) angle estimator; the resulting estimated angle is the one most likely to produce the intensity reading observed over the whole sweep (which might not happen to align with the angle of strongest signal if the other, less loud directions give contradictory information).

Once an angle is deduced from a sweep (in whatever way is previously chosen), there is further design freedom in how to use this estimate. A heuristic greedy strategy that circumvents the need for explicit target location estimation is to follow the estimated direction until a new sweep is completed and the angle is updated; we see it applied in the work of \citet{venkateswaran2013rf} who mount a directional
antenna on a naturally spinning micro aerial vehicle. They show that, given some bounds on the noise and frequency of outliers, the vehicle returns to a small ball around the source (whose radius grows with the distance traveled between each sweep) infinitely often. The strategy proves remarkably robust to the error of individual estimates: 23\% of
the bearings were off by more than $60^\circ$, yet a rover was able to reach a source initially 55\,m away in only 59.9\,m of flight, the large angle errors canceling out over the course of the run. 

The more involved way is to actually estimate source position and refine this estimate with each sweep. This strategy can sometimes interfere with robot motion: there is incentive in moving the robot to a ``good measurement spot", a spot that produces new information and allows the location estimate to converge quickly. This movement goal might not necessarily contribute to bringing the robot closer to the source, and there is a tradeoff between minimizing unnecessary movement on the path to the source and exploiting robot motion to get varied and useful measurement. \citet{feldbaum1960dual} calls this the dual effect of control, and optimal dual control is not generally tractable.

\citet{deghat2014localization} explicitly frame their estimator\textbackslash controller design task as dual control. In their setup, the observer must estimate the location of a target (which may drift slowly), then it must circle it at a given radius. Their estimator corrects the error perpendicular to the bearing while the controller moves closer to the target (sideways, in order to always rotate the bearing). They prove that the estimation error converges exponentially if and only if the bearing keeps sweeping through all directions (\emph{persistence of excitation}, a condition that circling the target guarantees).

Wildlife-tracking literature, where robots locate radio-tagged animals, provides many useful examples of dual control. \citet{tokekar2011active} use a ground robot equipped with a rotating directional antenna to locate invasive fish under lake ice, choosing sensing locations to minimize the
uncertainty of the estimate. 

\citet{vanderhook2014cautious} minimize
the total time spent traveling and measuring, landing on a ``cautious greedy'' strategy. Briefly, the target's position is estimated with an extended Kalman filter (EKF) whose uncertainty looks like an ellipse around the current estimate. At each step, the robot moves to a point
from which its line of sight to the estimate is perpendicular to the long axis of the uncertainty ellipse (this is the direction in which the target's position is least known), so that each new bearing reduces the largest uncertainty. It also stays far enough from the current estimate that the target
is unlikely to be behind it (to minimize the effect of the antenna's front--back ambiguity). Since this safe distance shrinks with the uncertainty, the robot effectively zigzags towards the targets as it locates it, meaning locating and approaching happen along the same path. The strategy is shown to be within a constant factor of the best achievable (time-wise). A cone with a single opening (i.e. without front--back ambiguity) performs even better, although reflections can still produce misleading
spurious peaks.
\citet{vanderhook2015algorithms} extend the analysis to several robots under communication constraints.

\citet{chen2023conservationbots} track several animals with an aerial robot and combine angle of arrival (AoA) and received signal strength indicator (RSSI) methods. They design a planner that jointly decides where to fly and which type of measurement to take (angle or strength), and note the effect of target motion on that choice: sensor rotations, which are slow to perform, are more accurate for static animals but become
counterproductive when the target moves.

The main takeaway here is that a single bearing only constrains the source to a line, and it is the sensor's motion between
measurements that pins down its exact position. The literature on submarine sonars can be surprisingly informative in determining what makes this motion ``useful" (from the point of view of maximizing information extracted from measurements). 

\citet{lingren1978position} analyzed whether the position and velocity of a target ship can be deduced by the bearing measurements of an observer ship, and \citet{nardone1981observability} gave
necessary and sufficient conditions on the observer's motion for a constant-velocity target. They conclude that a constant-velocity observer cannot always determine a constant-velocity target (acceleration is necessary), and that motion along the line of sight is not sufficient (even in the presence of observer acceleration). For a stationary source, two bearings that do not lie on the line of sight are sufficient. 
While the rover in our setup has access to intensity readings (unlike the ships in this analysis), the literature on submarine sonar supports the usefulness of occasionally leaving the line of sight path (the shortest path to the source) in order to gain confidence in the location estimate. This is reminiscent of the zigzag cautious greedy strategy of \citet{vanderhook2014cautious}.


\subsection{Estimating the Target Position}

The estimators relevant to our setup are approximations of the recursive Bayes filter, which, at each step, predicts the target position given its internal model of the problem dynamics then updates its belief with each new measurement depending on how much they align with the prediction (i.e. depending on the conditional probability of receiving the measurement given the current belief about the target location). \citet{thrun2005probabilistic}, \citet{barshalom2001estimation}, and \citet{simon2006optimal} derive various useful models for different estimation scenarios.

The main filters of interest mostly differ in the way they model belief about the target position: the Kalman filter \citep{kalman1960new} is for linear models; the Extended Kalman Filter (EKF) tolerates nonlinear models by linearizing them around the current estimate; the Unscented Kalman Filter (UKF) \citep{julier2004unscented} propagates a small set of sample points through the nonlinear model (rather than linearizing it); particle filters \citep{gordon1993novel,arulampalam2002tutorial} represent their belief about target location as weighted samples (each sample being a hypothesis), which makes them able to keep track of several strong candidates until the uncertainty is resolved (a relevant feature in the case of very strong reflections for instance); and grid filters break the map into cells and store for each cell the probability that the target is located in it.

\citet{aidala1983utilization} show that a Cartesian-coordinate EKF in a bearing-only setup can become overconfident and diverge, and propose the use of modified polar coordinates to separate observable (without the need for motion) quantities (like bearing) from the unobservable (requiring observer motion) ones (like range). \citet{ristic2004beyond} compare the EKF, the UKF, and particle filters on bearing-only tracking.


\subsection{Control of a Differential-Drive Rover}
The chosen observer is a differential-drive rover, and strategies for controlling it are well-established, usually via Model predictive control (MPC)
\citep{mayne2000constrained,rawlings2017model}. Different strategies specific to a moving target also exist (i.e. ways to reach it more efficiently than simply aiming for its current position) \citep{zarchan2012tactical,deghat2014localization,shames2012circumnavigation}. We do not expand much on this aspect of the control at the present time due to length constraints.

\subsection{Summary and Literature Gap}
Overall, there is significant work on Sound Source Localization, leveraging rotating sensors and robot motion for location estimation, and dual control. The dynamics of differential-drive rovers are also well understood (although skipped). The specific setup remains unique, especially given the choice of acoustic waves (as opposed to radio), and the many degrees of freedom in the design of the estimator and controller (the use of an estimator at all being itself a choice) mean that many specific combinations have not been validated.


